\documentclass[10pt,a4paper]{amsart}
\usepackage[margin=19mm]{geometry}
\usepackage{amsmath,amssymb,mathtools}
\usepackage[scr=rsfs,cal=euler]{mathalfa}
\usepackage{xfrac}
\usepackage[unicode,hidelinks,bookmarks=false]{hyperref}
\newcommand{\ie}{i.e.\ }
\newcommand{\cf}{cf.\ }
\newcommand{\resp}{respectively\ }
\newcommand{\Subsection}{\subsection}
\providecommand{\nobreakdash}{\nobreak\hbox{-}}
\setlength{\emergencystretch}{2em}
\multlinegap0pt
\usepackage{amssymb}
\usepackage{tikz}
\newtheorem{proposition}{Proposition}
\newcommand{\fr}{\mathfrak{fr}}
\newcommand{\frone}{\mathfrak{fr}^{>1}_k(x_0,x_1)}
\title{Polylogarithmic characterizations of the reduced coaction Lie algebra and the double shuffle Lie algebra, and their relation to the Kashiwara--Vergne Lie algebra}
\author{Megan Howarth, Muze Ren}
\date{Source: arXiv:2509.20275v3 (CC BY 4.0)}
\begin{document}
\maketitle

\noindent\emph{Source and licence.} Polylogarithmic characterizations of the reduced coaction Lie algebra and the double shuffle Lie algebra, and their relation to the Kashiwara--Vergne Lie algebra. Version arXiv:2509.20275v3 (CC BY 4.0), distributed by arXiv under the Creative Commons Attribution 4.0 International licence, http://creativecommons.org/licenses/by/4.0/, as recorded in the arXiv licence metadata for this exact version and shown on the abstract page https://arxiv.org/abs/2509.20275v3. Full licence notice: \texttt{licenses/arxiv-2509.20275-CC-BY-4.0.txt}.\par\smallskip

\noindent\textit{Editorial scope.} Counted unit: the article's proposition prop:Hex-Lie-subalgebra (source0.tex lines 2552--2576). The article's complete original proof of it is delivered with it (source0.tex lines 2578--2710, one \texttt{proof} environment of 133 lines). No mathematics is written, repaired or re-derived: the statement and the proof are the article's own lines, in the article's own notation, and no retained line is substituted or re-flowed.  Retained uncounted as the source-local context the proof needs: lines 508--510.  Nothing was omitted from any retained span, so the package carries no omission record.  No retained line contains a citation, so the wrapper carries no bibliography and no bibliography entry.  No retained line contains a figure environment, an image-inclusion command or drawing code, so no image is delivered and the package declares no asset dependency.  Editorial formatting only, in addition: an AMS \texttt{amsart} preamble that restates the article's own declarations and macros used by the retained lines, namely the environments \texttt{proposition} on separate counters, together with the standard packages those lines need.  The retained environments print in this document as Proposition 1 in order of appearance, whereas the article prints the same items with the numbering of its own full document.  The complete native source of the article as distributed in the arXiv e-print for this exact version is bundled unchanged as \texttt{sources/185673/source0.tex}.
	\begin{equation}\label{eq:Ihara_bracket}
		\{\psi_1,\psi_2\}=d_{\psi_2}(\psi_1)-d_{\psi_1}(\psi_2)-[\psi_1,\psi_2],
	\end{equation}

\begin{proposition}\label{prop:Hex-Lie-subalgebra}


The vector space
\[
\mathrm{Vep}
=
\left\{
\psi\in\mathfrak{fr}^{>1}_k(x_0,x_1)
\;\middle|\;
[x_1,\psi(x_\infty,x_1)]
+
[x_0,\psi(x_\infty,x_0)]
=0
\right\}
\]
is a Lie subalgebra of $\frone$ for the Ihara bracket \eqref{eq:Ihara_bracket}.
Moreover, the map
$L_1:
\mathrm{Vep}\to\mathfrak{sder}_2$ is an injective Lie algebra homomorphism, % which means 
i.e. for all \(\psi_1,\psi_2\in\mathrm{Vep}\),
\[
L_1(\{\psi_1,\psi_2\}) = [L_1(\psi_1),L_1(\psi_2)].
\]
\end{proposition}

\begin{proof}
We start by rewriting the Vep condition. Let 
\(\psi\in\mathfrak{fr}^{>1}_k(x_0,x_1)\), then
\begin{align*}
u_\psi(x_\infty)
&=
-u_\psi(x_0)-u_\psi(x_1)=
[x_0,\psi(x_\infty,x_0)]
+[x_1,\psi(x_\infty,x_1)],
\end{align*}
hence
\[
\psi\in\mathrm{Vep}
\quad\Longleftrightarrow\quad
u_\psi(x_\infty)=0
\quad\Longleftrightarrow\quad
u_\psi\in\mathfrak{sder}_2.
\]

Let \(\psi_1,\psi_2\in \mathrm{Vep}\), then \(u_{\psi_j}(x_\infty)=0\)
%and denote $u_j:=L_1(\psi_j)$, 
for $j \in \{1,2\}$ and it follows that for 
%Since \(u_j(x_\infty)=0\), for 
every
\(\varphi\in\mathfrak{fr}_k(x_0,x_1)\) and
\(i\in\{0,1\}\), we have
\begin{equation}\label{eq:substitution-derivation}
u_{\psi_j}\bigl(\varphi(x_\infty,x_i)\bigr)
=
-(d_{\psi_j}\varphi)(x_\infty,x_i),
\end{equation}
where $d_{\psi_j}$ is the derivation defined in \eqref{eq:Ihara_bracket}.
Indeed, after the substitution $x_0 \mapsto x_\infty, x_1 \mapsto x_i$ (for $i\in \{0,1\}$), the two sides of
\eqref{eq:substitution-derivation} are derivations which agree on the two free
generators: the first is sent to
\(u_{\psi_j}(x_\infty)=0\), while the second is sent to
$u_{\psi_j}(x_i)
= -[x_i,\psi_j(x_\infty,x_i)].$

For \(i\in\{0,1\}\), we compute
\begin{align*}
[u_{\psi_1},u_{\psi_2}](x_i)
={}&
- u_{\psi_1}\bigl([x_i,\psi_2(x_\infty,x_i)]\bigr)
+
u_{\psi_2}\bigl([x_i,\psi_1(x_\infty,x_i)]\bigr)                 \\
={}&
\bigl[[x_i,\psi_1(x_\infty,x_i)],
      \psi_2(x_\infty,x_i)\bigr]-
\bigl[[x_i,\psi_2(x_\infty,x_i)],
      \psi_1(x_\infty,x_i)\bigr]                      +
\left[
x_i,
\bigl(
d_{\psi_1}(\psi_2)-d_{\psi_2}(\psi_1)
\bigr)(x_\infty,x_i)
\right] \\
={}& \left[
x_i,
[\psi_1,\psi_2](x_\infty,x_i)
\right] + \left[
x_i,
\bigl(
d_{\psi_1}(\psi_2)-d_{\psi_2}(\psi_1)
\bigr)(x_\infty,x_i)
\right],
\end{align*}
where the third equality follows from the Jacobi identity.
% By the Jacobi identity,
% \begin{align*}
% &
% \bigl[[x_i,\psi_1(x_\infty,x_i)],
%       \psi_2(x_\infty,x_i)\bigr]-
% \bigl[[x_i,\psi_2(x_\infty,x_i)],
%       \psi_1(x_\infty,x_i)\bigr] \\
% &=
% \left[
% x_i,
% [\psi_1,\psi_2](x_\infty,x_i)
% \right].
% \end{align*}
Therefore,
\begin{align*}
[u_{\psi_1},u_{\psi_2}](x_i)
&=
\left[
x_i,
\bigl(
[\psi_1,\psi_2]
+d_{\psi_1}(\psi_2)
-d_{\psi_2}(\psi_1)
\bigr)(x_\infty,x_i)
\right]                                                     \\
&=
\left[
x_i,
-\{\psi_1,\psi_2\}(x_\infty,x_i)
\right],
\end{align*}
by definition of the Ihara bracket.
% \[
% \{\psi_1,\psi_2\}
% =
% d_{\psi_2}(\psi_1)
% -d_{\psi_1}(\psi_2)
% -[\psi_1,\psi_2].
% \]
Hence  $\rho([L_1(\psi_1), L_1(\psi_2)]) = \rho(L_1(\{\psi_1, \psi_2\}))$, and since $\rho$ is injective on $L_1(\fr_k^{>1})$, we conclude $[L_1(\psi_1),L_1(\psi_2)]$ %$[u_{\psi_1},u_{\psi_2}]
$=L_1(\{\psi_1,\psi_2\}).$ 
Moreover, it follows from the previous computations that
\begin{align*}
    u_{\{\psi_1, \psi_2\}}(x_\infty) = [u_{\psi_1},u_{\psi_2}](x_\infty) = u_{\psi_1}(u_{\psi_2}(x_\infty)) - u_{\psi_2}(u_{\psi_1}(x_\infty)) = 0,
\end{align*}

% Since \(u_1,u_2\in\mathfrak{sder}_2\) and
% \(\mathfrak{sder}_2\) is a Lie subalgebra of
% \(\mathfrak{tder}_2\), their commutator fixes \(x_\infty\),
% $[u_1,u_2](x_\infty)=0.$

% The previous identity therefore gives $L_1(\{\psi_1,\psi_2\})(x_\infty)=0$,
which is equivalent to
$\{\psi_1,\psi_2\}\in\mathrm{Vep}.$ Thus, \(\mathrm{Vep}\) is a Lie subalgebra %and $[L_1(\psi_1),L_1(\psi_2)] = L_1(\{\psi_1,\psi_2\})$,
and $L_1$ is a Lie algebra homomorphism. Finally, it is injective because the change of variables
\[
\psi(x_0,x_1)
\mapsto
\psi(-x_0-x_1,x_1)
=
\psi(x_\infty,x_1)
\]
is an involutive automorphism of
\(\mathfrak{fr}_k(x_0,x_1)\).
\end{proof}

\end{document}
